\documentclass[11pt]{article}
\usepackage{amsmath, amssymb}
\usepackage{geometry}
\geometry{margin=1in}

\begin{document}

\begin{center}
\textbf{CSE400 – Fundamentals of Probability in Computing} \\[0.2cm]
\textbf{Lecture 6: Discrete Random Variables, Expectation and Problem Solving} \\
\textbf{Instructor:} Dhaval Patel, PhD \\
\textbf{Date:} January 22, 2025
\end{center}

\hrule
\vspace{0.4cm}

\section*{1. Random Variables}

\subsection*{Definition}

A \textbf{random variable} $(X)$ on a sample space $(\Omega)$ is a function
\[
X : \Omega \rightarrow \mathbb{R}
\]
that assigns to each sample point $(\omega \in \Omega)$ a real number $(X(\omega))$.

\subsection*{Discrete Random Variables}

In this lecture, attention is restricted to \textbf{discrete random variables}, i.e., random variables that take values in a range that is \textbf{finite or countably infinite}.

Although $(X)$ is defined as a mapping into $(\mathbb{R})$, the actual set of values
\[
\{ X(\omega) : \omega \in \Omega \}
\]
is a discrete subset of $(\mathbb{R})$.

\subsection*{Visualization of a Distribution}

\begin{itemize}
\item The x-axis represents the values that the random variable can take.
\item The height of the bar at a value $(a)$ is the probability $(\Pr[X = a])$.
\item Each probability is computed as the probability of the corresponding event in the sample space.
\end{itemize}

\section*{2. Discrete and Continuous Random Variables}

\subsection*{Discrete Random Variable}

\begin{itemize}
\item Countable support
\item Probability Mass Function (PMF)
\item Probabilities are assigned to single values
\item Each possible value has strictly positive probability
\end{itemize}

\subsection*{Continuous Random Variable}

\begin{itemize}
\item Uncountable support
\item Probability Density Function (PDF)
\item Probabilities are assigned to intervals of values
\item Each individual value has probability zero
\end{itemize}

\section*{3. Example: Tossing Three Fair Coins}

\subsection*{Experiment}

Three fair coins are tossed. \\
Let $(Y)$ denote the number of heads that appear.

\subsection*{Possible Values}

\[
Y \in \{0, 1, 2, 3\}
\]

\subsection*{Computation of Probabilities}

Each outcome of three coin tosses has probability $\left(\frac{1}{2}\right)^3 = \frac{1}{8}$.

\begin{itemize}
\item $\Pr[Y = 0] = \Pr[(t,t,t)] = \frac{1}{8}$
\item $\Pr[Y = 1] = \Pr[(t,t,h),(t,h,t),(h,t,t)] = \frac{3}{8}$
\item $\Pr[Y = 2] = \Pr[(t,h,h),(h,t,h),(h,h,t)] = \frac{3}{8}$
\item $\Pr[Y = 3] = \Pr[(h,h,h)] = \frac{1}{8}$
\end{itemize}

\subsection*{Legitimacy of the PMF}

Since $(Y)$ must take one of the values $(0,1,2,3)$,
\[
\sum_{i=0}^{3} \Pr[Y=i] = 1.
\]

\section*{4. Probability Mass Function (PMF)}

\subsection*{Definition}

A random variable that can take on at most a countable number of possible values is said to be \textbf{discrete}.

Let $(X)$ be a discrete random variable with range
\[
R_X = \{x_1, x_2, x_3, \dots\} \quad \text{(finite or countably infinite)}.
\]

The function
\[
P_X(x_k) = \Pr[X = x_k], \quad k = 1,2,3,\dots
\]
is called the \textbf{Probability Mass Function (PMF)} of $(X)$.

\subsection*{Property}

Since $(X)$ must take exactly one of the values in its range,
\[
\sum_{k=1}^{\infty} P_X(x_k) = 1.
\]

\section*{5. PMF Example with Parameter $(\lambda)$}

\subsection*{Given}

The PMF of a random variable $(X)$ is
\[
p(i) = c \frac{\lambda^i}{i!}, \quad i = 0,1,2,\dots
\]
where $(\lambda > 0)$.

\subsection*{Determination of the Constant $(c)$}

Since probabilities must sum to 1,
\[
\sum_{i=0}^{\infty} p(i) = 1.
\]
Substituting the PMF,
\[
c \sum_{i=0}^{\infty} \frac{\lambda^i}{i!} = 1.
\]
Using the identity
\[
e^{\lambda} = \sum_{i=0}^{\infty} \frac{\lambda^i}{i!},
\]
we obtain
\[
c e^{\lambda} = 1 \quad \Rightarrow \quad c = e^{-\lambda}.
\]

\subsection*{Computation of Probabilities}

\[
\Pr[X = 0] = e^{-\lambda} \frac{\lambda^0}{0!} = e^{-\lambda}.
\]

\[
\Pr[X > 2] = 1 - \Pr[X \le 2]
\]
\[
= 1 - \big(\Pr[X=0] + \Pr[X=1] + \Pr[X=2]\big)
\]
\[
= 1 - \left(e^{-\lambda} + \lambda e^{-\lambda} + \frac{\lambda^2}{2} e^{-\lambda}\right).
\]

\section*{6. Bayes’ Theorem}

\subsection*{Statement}

Using
\[
\Pr(A \cap B) = \Pr(B \mid A)\Pr(A),
\]
we obtain \textbf{Bayes’ Formula}:
\[
\Pr(B_i \mid A) =
\frac{\Pr(A \mid B_i)\Pr(B_i)}
{\sum_j \Pr(A \mid B_j)\Pr(B_j)}.
\]

\subsection*{Terminology}

\begin{itemize}
\item $(\Pr(B_i))$: \emph{a priori probability}
\item $(\Pr(B_i \mid A))$: \emph{posterior probability}
\end{itemize}

\section*{7. Bayes’ Theorem Example: Auditorium with 30 Rows}

\subsection*{Description}

\begin{itemize}
\item Auditorium has 30 rows.
\item Row 1 has 11 seats, Row 2 has 12 seats, \dots, Row 30 has 40 seats.
\item A door prize is awarded by:
\begin{enumerate}
\item Randomly selecting a row (each row equally likely),
\item Randomly selecting a seat within that row (each seat equally likely).
\end{enumerate}
\end{itemize}

\subsection*{Computations}

\textbf{Probability of Seat 15 given Row 20}

Row 20 has 30 seats, hence
\[
\Pr(\text{Seat 15} \mid \text{Row 20}) = \frac{1}{30}.
\]

\textbf{Probability of Row 20 given Seat 15}

Using Bayes’ formula,
\[
\Pr(\text{Row 20} \mid \text{Seat 15})
= \frac{\Pr(\text{Seat 15} \mid \text{Row 20}) \Pr(\text{Row 20})}
{\Pr(\text{Seat 15})}.
\]

From the lecture computation,
\[
\Pr(\text{Seat 15}) \approx 0.0342,
\quad
\Pr(\text{Row 20} \mid \text{Seat 15}) \approx 0.0325.
\]

\end{document}

